\clearpage
\section*{Sources and further reading}
\addcontentsline{toc}{section}{Sources and further reading}

The proofs above use the displayed coefficient calculations. These references provide related examples and the wider background.

\begin{enumerate}[leftmargin=*,label={[\arabic*]},itemsep=.65em]
\item Henry Shin, \href{https://refactor.press/refactors/coordinate-repair}{\emph{How far can a change of coordinates go?}}, Refactor. The starting example \(x+y+xy\) in arithmetic modulo three. The companion \href{https://refactor.press/refactors/first-order}{\emph{Make first order exact}} develops the practice of retaining specified coefficients. The longer \href{https://refactor.press/notes/when-a-formal-logarithm-fails}{\emph{When a Formal Logarithm Fails}} follows the original coordinate-repair calculation further.

\item Charles Rezk, \href{https://www.sas.rochester.edu/mth/sites/doug-ravenel/otherpapers/hopkins-miller-thm.pdf}{\emph{Notes on the Hopkins--Miller Theorem}}, \S3.1, printed p.~3, for coordinate changes of addition rules; \S3.5, printed pp.~3--4, for the information obtained by repeating an operation \(p\) times; \S5.1--5.3, printed pp.~7--8, for corrections with a square-zero coefficient and the polynomial patterns used in the degree-nine example.

\item Ariane M\'ezard, Matthieu Romagny, and Dajano Tossici, \href{https://www.numdam.org/item/10.5802/jtnb.863.pdf}{\emph{Sekiguchi--Suwa theory revisited}}, Journal de Th\'eorie des Nombres de Bordeaux \textbf{26} (2014), 163--200, \S2.2, printed p.~167. Background on the family connecting addition with multiplication through a coefficient parameter. The exact repair criterion used here is proved in Section~\ref{sec:exact-theorem}.

\item Jacob Lurie, \href{https://www.math.ias.edu/~lurie/252xnotes/Lecture13.pdf}{\emph{The Stratification of MFG}}, Lecture 13, April 27, 2010, printed p.~2, Remark~4. The comparison of leading coefficients after repeating an operation is the general background to the direct calculation \(\lambda'=\lambda/c^2\) in Sections~\ref{sec:classification}--\ref{sec:three-series}.
\end{enumerate}
